% page 553 du fac-simile (image p-552 du scan)
% bloc 162.6 x 242.0 mm ; marge gauche 39.4 mm
\pgc{17.780mm}{0.000mm}{
\l{\cel{52}545}
\l{}
\l{\sou{spekular}. (\sou{netrans., pri}).I. Agar explori teoriala. - II.}
\l{\cel{3}(\sou{borso}).Facar operaci (kompri, vendi) riskoza, fundamento di}
\l{\cel{3}chanci di kresko o di deskresko di la varo-preci, o di la}
\l{\cel{3}aciono-preci. - DEFIRS.}
\l{}
\l{\sou{spelto}. (\sou{bot.}) Speco di frumento di qua la brano duras adherar}
\l{\cel{3}a la grano pos drasheso. - L.\cel{1}\sou{triticum spelta}. DEIS.}
\l{}
\l{\sou{spencero}. I. Frako sen baski. - II. Korsajo sen jupo. DEFI.}
\l{}
\l{\sou{spensar}. (\sou{tra}ns.)I.Donar, uzar(pekunio) por prokurar a su ulo.}
\l{\cel{3}- II. (\sou{metaf}.)Donar, uzar (sua tempo, esforci, fakultati)}
\l{\cel{3}por obtenar la rezultajo quan onu deziras. - EFI.}
\l{}
\l{\sou{spergulo}. (\sou{bot.}) Planto ek la familio \textquotedbl{}kariofilei\textquotedbl{}. - L.\cel{1}\sou{sper}-}
\l{\cel{3}\sou{gula}. DEFIS.}
\l{}
\l{\sou{spermo}. (\sou{fiziol.}) Liquoro fekundigiva de la maskulo. - DEFIRS.}
\l{}
\l{\sou{spermaceto}. Materio perlomatrea, extraktita de oleo qua bundegas}
\l{\cel{3}en la kavaji di la kapo di la kashaloto, ed uzata por facar}
\l{\cel{3}bujii\cel{2}luxoza, koldkremo, e c. - DEFIR.}
\l{}
\l{\sou{spermatido}. (\sou{biol.}) Spermatozoido nova, derivita de la duopla}
\l{\cel{3}du-partesko inter-sucedanta di la spermatociti. -}
\l{}
\l{\sou{spermatocito}. (\sou{biol.}) Spermatogonio qua, cesante partigar su,}
\l{\cel{3}komencas kreskar e donas, pose, per du duparteski intersuce-}
\l{\cel{3}danta, la spermatidi. - DEFIS.}
\l{}
\l{\sou{spermatoforo}. (\sou{biol}.) Kapsulo qua kontenas pakaji de spermato-}
\l{\cel{3}zoidi e qua renkontresas precipue che la cefalopodi. - DEFIS.}
\l{}
\l{\sou{spermatogenezo}. (\sou{biol.}) Ensemblo ek la procesi evolucionala qua}
\l{\cel{3}abutas a la formaco di la spermatozoidi. - DEFIS.}
\l{}
\l{\sou{spermatogonio.}\cel{2}(\sou{biol.}) Celulo di la tubo seminifera qua su}
\l{\cel{3}partigas multega-foye e deskreskas pri sua volumino : lu pro-}
\l{\cel{3}duktas pose la spermatociti. DEFIS.}
\l{}
\l{\sou{spermatoreo}.(patol.) Emisi spermala ne-volata. - DEFIS.}
\l{}
\l{\sou{spermatozoido}. (\sou{biol}.)Animaleto quan onu trovas en la spermo}
\l{\cel{3}di la animali, ed en la poleno qua fekundigas la planti.}
\l{\cel{3}- DEFIS.}
\l{}
\l{\sou{spermino}. (\sou{farmacio)}. Bazo ne bone analizita qua esus la subs-}
\l{\cel{3}tanco agiva di la preparaji opoterapiala quin\cel{1}\sou{onu}\cel{1}obtenas ek}
\l{\cel{3}la testikli, ovarii, spleno, prostato, e c. - DEFIS.}
\l{}
\l{spermocentro. (\sou{biol.}) Centrosomo di spermatozoido, pri qua}
\l{\cel{3}Plättner dicis ke lu esas renkontrebla en la pinto antea di}
\l{\cel{3}la kapo, kontraste kun la dico da Flemming e Hertwig, segun}
\l{\cel{3}qui lu esus en la segmento meza. -}
}
